\documentclass[10pt,a4paper]{amsart}
\usepackage[margin=19mm]{geometry}
\usepackage{amsmath,amssymb,mathtools}
\usepackage[scr=rsfs,cal=euler]{mathalfa}
\usepackage{xfrac}
\usepackage[unicode,hidelinks,bookmarks=false]{hyperref}
\newcommand{\ie}{i.e.\ }
\newcommand{\cf}{cf.\ }
\newcommand{\resp}{respectively\ }
\newcommand{\Subsection}{\subsection}
\providecommand{\nobreakdash}{\nobreak\hbox{-}}
\setlength{\emergencystretch}{2em}
\multlinegap0pt
\usepackage{bm}
\usepackage{bbm}
\usepackage{dsfont}
\usepackage{xspace}
\usepackage{cancel}
\usepackage{mathtools}
\usepackage{amsthm}
\makeatletter
\@ifundefined{Theorem}{\newtheorem{Theorem}{Theorem}}{}
\@ifundefined{Lem}{\newtheorem{Lem}[Theorem]{Lemma}}{}
\makeatother
\newcommand{\G}{\mathbb{G}}
\newcommand{\Hh}{\mathbb{H}}
\newcommand{\Hsp}{\mathcal{H}}
\newcommand{\Ind}{\mathrm{Ind}}
\newcommand{\X}{\mathbb{X}}
\newcommand{\id}{\mathrm{id}}
\newcommand{\msG}{\mathscr{G}}
\newcommand{\ov}{\overline}
\title[The standard construction for cocycle twisted and braided te]{The standard construction for cocycle twisted and braided tensor product W$^*$-algebras}
\author{K. De Commer, J. Krajczok}
\date{Source: arXiv:2508.00595v1 (CC BY 4.0)}
\begin{document}
\maketitle

\noindent\emph{Source and licence.} The standard construction for cocycle twisted and braided tensor product W$^*$-algebras. Version arXiv:2508.00595v1 (CC BY 4.0), distributed under the Creative Commons Attribution 4.0 International licence, as recorded in the arXiv licence metadata for this exact version and shown on the abstract page https://arxiv.org/abs/2508.00595v1. Full rights notice: licenses/arxiv-2508.00595-CC-BY-4.0.txt.\par\smallskip

\noindent\textit{Editorial scope.} Counted unit: the Lem whose source label is \texttt{LemIndLModConj} in the article (source0.tex lines 1521--1536, source label \texttt{LemIndLModConj}), delivered together with the article's own complete original proof of it (source0.tex lines 1538--1625, one \texttt{proof} environment of 88 lines). No mathematics is written, repaired or re-derived: statement and proof are the article's own text line for line. Required context retained uncounted: EqStandardMorita (lines 536--544); EqMoritaModCon (lines 554--556); EqSplitMatrixCross (lines 936--939); EqIsometImpl (lines 987--989); eq50 (lines 1078--1080); TheoIndL (lines 1433--1442). Every definition, display and label of those lines is retained unchanged. Omitted from this package and not redistributed: 1--535, 545--553, 557--935, 940--986, 990--1077, 1081--1432, 1443--1520, 1537--1537, 1626--4309. Nothing inside a retained excerpt is omitted or altered: every retained body is delivered exactly as the article's lines are written, so no omission receipt of any kind accompanies this package. Citations: the retained excerpts carry no citation command at all, so no bibliography entry of the article is transcribed and no reference is redistributed. Reference adaptation: none was needed and none was made; every reference inside the delivered unit resolves inside it, so no unresolved reference or citation is produced. Printed slips kept literal: no printed word, symbol, hypothesis or number of the counted statement or of its proof was altered. Layout and class: the article's own class is \texttt{amsart} with packages including mathabx, tikz-cd, verbatim, fullpage, dsfont, hyperref, parskip, amsrefs, pdfpages, multirow, array, mystix2, amssymb, amsfonts; the wrapper is the lane's pinned \texttt{amsart} editorial wrapper at 10pt on A4, which restates the theorem-like environments the retained text needs and adds the article's own macro definitions that the retained text needs (\texttt{\textbackslash G}, \texttt{\textbackslash Hh}, \texttt{\textbackslash Hsp}, \texttt{\textbackslash Ind}, \texttt{\textbackslash X}, \texttt{\textbackslash id}, \texttt{\textbackslash msG}, \texttt{\textbackslash ov}), with extra wrapper packages (bm, bbm, dsfont, xspace, cancel, mathtools). The delivered numbering is the unit's own. Assets: no retained span contains a figure environment, a graphics-inclusion command, a PDF-page inclusion, a table or a TikZ picture; no image file is copied into this package, the package declares no asset dependency and no image file is redistributed with it. Licence: version arXiv:2508.00595v1 of this article is distributed under the Creative Commons Attribution 4.0 International licence, as recorded in the arXiv licence metadata for this exact version and shown on the abstract page https://arxiv.org/abs/2508.00595v1. A full rights notice is delivered at licenses/arxiv-2508.00595-CC-BY-4.0.txt. Characters: every retained excerpt is pure ASCII, so the delivered engine, XeLaTeX with its default Latin Modern text font, reports no missing character.
\begin{equation}\label{EqStandardMorita}
 \begin{pmatrix} L^2(\hat{P}) & L^2(\hat{N}) \\ L^2(\hat{O}) & L^2(\hat{M}) \end{pmatrix}
 \cong 
\begin{pmatrix} \hat{N} \bar{\otimes}_{\hat{M}}\overline{\Hsp}  & \Hsp \\ \overline{\Hsp} & L^2(\hat{M}) \end{pmatrix},
\qquad  
 \begin{pmatrix} \pi_{\hat{Q}}(x)\hat{J}_{12}\pi_{\hat{Q}}(w)\xi & \pi_{\hat{Q}}(y)\eta \\ \hat{J}_{12} \pi_{\hat{Q}}(z) \zeta & \vartheta\end{pmatrix}
 \mapsto
 \begin{pmatrix} x\otimes \overline{w\xi} & y\eta \\ \overline{z\zeta} & \vartheta\end{pmatrix},
\end{equation}

\begin{equation}\label{EqMoritaModCon}
J_{\hat{Q}}\begin{pmatrix} x\otimes \overline{y\xi} & \eta \\ \overline{\zeta} & \vartheta\end{pmatrix} = \begin{pmatrix} y\otimes \overline{xJ_{\hat{M}}\xi} & \zeta \\ \overline{\eta} & J_{\hat{M}}\vartheta\end{pmatrix}
\end{equation}

\begin{equation}\label{EqSplitMatrixCross}
\msG \ltimes A = \begin{pmatrix} \G_{11}\ltimes A_1 & \G_{12} \ltimes A_2 \\ 
\G_{21}\ltimes A_1 & \G_{22}\ltimes A_2\end{pmatrix},  \qquad \G_{ij}\ltimes A_j  = \hat{p}_i (\msG\ltimes A)\hat{p}_j. 
\end{equation}

\begin{equation}\label{EqIsometImpl}
U_{\gamma} = (J_{\hat{Q}}\otimes J_A)J_{\msG \ltimes A},
\end{equation}

\begin{equation}\label{eq50}
\gamma_2^1 =(\Delta_{22}^1\otimes \id) \gamma\colon A_2 \rightarrow L^{\infty}(\G_{21})\bar{\otimes} A_1,
\end{equation}

\begin{Theorem}\label{TheoIndL}
There is a canonical identification of unitary $\Hh$-representations
\[
I\colon (L^2(\Ind_{\X}(A)),U_{\Ind_{\X}(\gamma)}) \cong (\Ind_{\X}(L^2(A)),\Ind_{\X}(U_{\gamma})).
\]
Moreover, the induced left $\Ind_{\X}(A)$-representation on $\Ind_{\X}(L^2(A))$ is given by left multiplication;
\[
I\pi_{\Ind_{\X}(A)}(w)I^* \eta = (\id\otimes \pi_A)(w)\eta,\qquad \forall w\in \Ind_{\X}(A), \eta\in \Ind_{\X}(L^2(A)).  
\]
\end{Theorem}

\begin{Lem}\label{LemIndLModConj}
Let $(A,\gamma)$ be a $\G$-W$^*$-algebra. Under the identification in Theorem \ref{TheoIndL}, the modular conjugation $J_{\Ind_{\X}(A)}$ of $\Ind_{\X}(A)$ is the unique anti-linear map $\Ind_{\X}(L^2(A))\rightarrow \Ind_{\X}(L^2(A))$ such that 
\begin{multline}\label{EqFormModConj}
(\hat{J}_{11}\otimes J_{\Ind_{\X}(A)})
\bigl(
\hat{\pi}_{12}^1(x)_{[1]}
(\Delta_{22}^1\otimes \id)(\gamma(a)U_{\gamma}^*)_{[123]}
((\hat{J}_{12}\otimes J_A)(\hat{\pi}_{12}^2(z)\otimes 1)\gamma(b)\xi)_{[13]}\bigr) \\
=
\Ind_{\X}(U_{\gamma})
\hat{\pi}_{12}^1(z)_{[1]}
(\Delta_{22}^1\otimes \id)(\gamma(b)U_{\gamma}^*)_{[123]}
((\hat{J}_{12}\otimes J_A)(\hat{\pi}_{12}^2(x)\otimes 1)\gamma(a)U_\gamma^* (\hat{J}_{22}\otimes J_A)\xi)_{[13]}
\end{multline}
for $\xi\in L^2(\G)\otimes L^2(A)$, $x,z\in \hat{N}$ and $a,b\in A$. 
\end{Lem}

\begin{proof}
Applying \eqref{EqStandardMorita} to the W$^*$-Morita equivalence \eqref{EqSplitMatrixCross}, we obtain a composition of unitary identifications 
\begin{equation}\begin{split}\label{eq49}
(\G_{12}\ltimes A) \overline{\otimes}_{\G_{22}\ltimes A} \overline{L^2(\G_{12})\otimes L^2(A)} &\cong (\G_{12}\ltimes A) \overline{\otimes}_{\G_{22}\ltimes A} \overline{L^2(\G_{12}\ltimes A)}
\cong L^2(\G_{11}\ltimes \Ind_{\X}(A)) \\
&\cong L^2(\G_{11})\otimes L^2(\Ind_{\X}(A)) \cong L^2(\G_{11}) \otimes \Ind_{\X}( L^2(A)).
\end{split}\end{equation}
Following the precise prescription of the maps in \eqref{eq49}, we see that this identification becomes
\[\begin{split}
&\quad\;
\hat{p}_1(\pi_{\hat{Q}}(x)\otimes 1)
\Ind_{\msG}(\gamma)(a)\hat{p}_2
 \otimes \ov{
 \hat{p}_1(\pi_{\hat{Q}}(z)\otimes 1)
\Ind_{\msG}(\gamma)(b)\hat{p}_2 \xi
}\\
&\mapsto 
\hat{p}_1(\pi_{\hat{Q}}(x)\otimes 1)
\Ind_{\msG}(\gamma)(a)\hat{p}_2
J_{12}^{\ltimes} 
 \hat{p}_1(\pi_{\hat{Q}}(z)\otimes 1)
\Ind_{\msG}(\gamma)(b)\hat{p}_2 \xi\\
&\mapsto 
(1\otimes U^*_{\Ind_{\msG}(\gamma),12})
\bigl(
\hat{p}_1(\pi_{\hat{Q}}(x)\otimes 1)
\Ind_{\msG}(\gamma)(a)\hat{p}_2
J_{12}^{\ltimes} 
 \hat{p}_1(\pi_{\hat{Q}}(z)\otimes 1)
\Ind_{\msG}(\gamma)(b)\hat{p}_2 \xi
\bigr)_{[13]}
\end{split}\]
for $x,z\in \hat{N},a,b\in A, \xi\in L^2(\G)\otimes L^2(A)$ (the first and the third isomorphisms in \eqref{eq49} are formal identities). Let us remark that we consider $L^2(\Ind_{\X}(A))$ as having one leg in this computation. Using \eqref{EqIsometImpl}, and by tracing to which space particular vectors belong to, we can simplify this expression to
\[\begin{split}
&\quad\;
U^*_{\Ind_{\msG}(\gamma),12 [23]}
\pi_{\hat{Q}}(x)_{[1]}
\Ind_{\msG}(\gamma)_2^1 (a)_{[13]}
U^*_{\Ind_{\msG}(\gamma),21[13]}
\bigl(
(\hat{J}_{12}\otimes J_A)
(\pi_{\hat{Q}}(z)\otimes 1)
\gamma(b) \xi
\bigr)_{[13]}\\
&=
\pi_{\hat{Q}}(x)_{[1]}
(\id\otimes \Ind_{\msG}(\gamma)_1^2)\Ind_{\msG}(\gamma)_2^1 (a)
U^*_{\Ind_{\msG}(\gamma),12 [23]}
U^*_{\Ind_{\msG}(\gamma),21[13]}
\bigl(
(\hat{J}_{12}\otimes J_A)
(\pi_{\hat{Q}}(z)\otimes 1)
\gamma(b) \xi
\bigr)_{[13]}.
\end{split}\]
Next, using that $\Ind_{\msG}(\gamma)^2_1$ is simply the inclusion map $\Ind_{\X}(A)\rightarrow L^{\infty}(\G_{12})\bar\otimes A$, we further obtain from \eqref{eq50} that
\[\begin{split}
&\quad\;
\pi_{\hat{Q}}(x)_{[1]}
\Ind_{\msG}(\gamma)_2^1 (a)
U^*_{\Ind_{\msG}(\gamma),12 [23]}
U^*_{\Ind_{\msG}(\gamma),21[13]}
\bigl(
(\hat{J}_{12}\otimes J_A)
(\pi_{\hat{Q}}(z)\otimes 1)
\gamma(b) \xi
\bigr)_{[13]}\\
&=
\pi_{\hat{Q}}(x)_{[1]}
(\Delta^1_{22}\otimes \id)\gamma(a)
(\Delta_{22}^1\otimes \id)(U^*_{\Ind_{\msG}(\gamma),22})
\bigl(
(\hat{J}_{12}\otimes J_A)
(\pi_{\hat{Q}}(z)\otimes 1)
\gamma(b) \xi
\bigr)_{[13]}\\
&=
\pi_{\hat{Q}}(x)_{[1]}
(\Delta^1_{22}\otimes \id)(\gamma(a)
U^*_{\gamma})
\bigl(
(\hat{J}_{12}\otimes J_A)
(\pi_{\hat{Q}}(z)\otimes 1)
\gamma(b) \xi
\bigr)_{[13]}.
\end{split}\]
Using \eqref{EqIsometImpl} and Theorem \ref{TheoIndL}, we see that under the isomorphisms \eqref{eq49}, the modular conjugation \eqref{EqMoritaModCon} is mapped to $\Ind_{\X}(U_{\gamma})^*(\hat{J}_{11}\otimes J_{\Ind_{\X}(A)})$. This gives \eqref{EqFormModConj}.
\end{proof}

\end{document}
